\documentclass[aspectratio=169,11pt]{beamer}

\usepackage[T1]{fontenc}
\usepackage{lmodern}
\usepackage{amsmath,amssymb,mathtools,bm}
\usepackage{booktabs,tabularx,array}
\usepackage{graphicx}
\usepackage{microtype}

\definecolor{SVEBlue}{HTML}{236BA6}
\definecolor{SVEGreen}{HTML}{207C61}
\definecolor{SVEOrange}{HTML}{CB7623}
\definecolor{SVEInk}{HTML}{20252B}
\definecolor{SVEMuted}{HTML}{5F6872}
\definecolor{SVEWash}{HTML}{F5F7F8}
\definecolor{SVERed}{HTML}{A33D3D}

\setbeamercolor{normal text}{fg=SVEInk,bg=white}
\setbeamercolor{structure}{fg=SVEBlue}
\setbeamercolor{frametitle}{fg=SVEInk,bg=white}
\setbeamercolor{title}{fg=SVEInk}
\setbeamercolor{subtitle}{fg=SVEMuted}
\setbeamercolor{block title}{fg=white,bg=SVEBlue}
\setbeamercolor{block body}{fg=SVEInk,bg=SVEWash}
\setbeamercolor{alerted text}{fg=SVERed}
\setbeamercolor{footline}{fg=SVEMuted,bg=white}
\setbeamerfont{title}{size=\fontsize{27}{31}\selectfont,series=\bfseries}
\setbeamerfont{subtitle}{size=\fontsize{14}{17}\selectfont}
\setbeamerfont{frametitle}{size=\fontsize{20}{23}\selectfont,series=\bfseries}
\setbeamerfont{normal text}{size=\fontsize{11}{13.3}\selectfont}
\setbeamerfont{footline}{size=\fontsize{7.5}{9}\selectfont}
\setbeamertemplate{navigation symbols}{}
\setbeamertemplate{itemize item}{\color{SVEBlue}$\blacktriangleright$}
\setbeamertemplate{itemize subitem}{\color{SVEGreen}$\vcenter{\hbox{\scriptsize$\bullet$}}$}
\setbeamertemplate{blocks}[rounded][shadow=false]
\setbeamertemplate{frametitle}{%
  \vspace{0.25cm}\insertframetitle\par
  \vspace{0.05cm}{\color{SVEBlue}\rule{\textwidth}{0.55pt}}}
\setbeamertemplate{footline}{%
  \leavevmode
  \hbox{%
    \begin{beamercolorbox}[wd=.78\paperwidth,ht=2.6ex,dp=1.2ex,leftskip=.55cm]{footline}%
      \insertshorttitle\hspace{0.8em}\textcolor{SVEBlue}{/}\hspace{0.8em}\insertsectionhead
    \end{beamercolorbox}%
    \begin{beamercolorbox}[wd=.22\paperwidth,ht=2.6ex,dp=1.2ex,rightskip=.55cm plus1fil]{footline}%
      \hfill\ifnum\value{framenumber}<18\insertframenumber\,/\,17\else Backup\,\number\numexpr\value{framenumber}-17\relax\fi
    \end{beamercolorbox}%
  }}

\newcommand{\E}{\mathcal E}
\newcommand{\pr}{\operatorname{pr}}
\newcommand{\D}{\mathrm D}
\newcommand{\Exact}{\textcolor{SVEGreen}{\bfseries Exact}}
\newcommand{\Numerical}{\textcolor{SVEBlue}{\bfseries Numerical}}
\newcommand{\Proposed}{\textcolor{SVEOrange}{\bfseries Proposed}}
\newcommand{\Printed}{\textcolor{SVERed}{\bfseries Printed source}}
\newcommand{\scope}[1]{\par\smallskip{\scriptsize\color{SVEMuted}\textbf{Scope.} #1}}
\newcommand{\source}[1]{\par\smallskip{\tiny\color{SVEMuted}#1}}

\title[Certified scalar SVE reductions]{Certified Lie-Point Reductions of a Scalar SVE Benchmark}
\subtitle{Exact symmetry classification, corrected profile equations,\\ local verification, and an ML research gate}
\author{Math 4890 dissertation demonstration}
\date{October 2026}

\begin{document}

\section{Dissertation}
\begin{frame}[plain]
  \vspace{0.65cm}
  {\color{SVEBlue}\rule{2.1cm}{1.2pt}}\par\vspace{0.72cm}
  {\usebeamerfont{title}\usebeamercolor[fg]{title}\inserttitle\par}
  \vspace{0.28cm}
  {\usebeamerfont{subtitle}\usebeamercolor[fg]{subtitle}\insertsubtitle\par}
  \vspace{0.8cm}
  {\small\insertauthor\par}\vspace{0.14cm}
  {\small\color{SVEMuted}\insertdate\par}
  \vfill
  {\small\color{SVEMuted}
  Printed scalar model $\longrightarrow$ unrestricted point symmetries
  $\longrightarrow$ invariant ODEs $\longrightarrow$ checked local solutions}
  \vspace{0.22cm}
\end{frame}
\begin{frame}{Research question and proof chain}
\small
\textbf{Question.} Which mathematically certified reductions of the \emph{printed scalar} Saint--Venant--Exner (SVE) model yield useful lower-dimensional solution problems?
\vspace{0.23cm}
\begin{center}
\setlength{\fboxsep}{8pt}
\fcolorbox{SVEBlue}{SVEWash}{\parbox{.89\linewidth}{\centering
$E_m=0\ \xrightarrow{\ \operatorname{pr}^{(3)}X\ }\ \mathfrak g_m$
$\ \xrightarrow{\ Q=0\ }\ \phi=\mathcal T[F]$\\[4pt]
$\xrightarrow{\ E_m[\mathcal TF]\ }\ \mathcal R[F]=0\ \xrightarrow{\ \text{verified integral}\ }\ I[F]=I_0$}}
\end{center}
\begin{columns}[T,onlytextwidth]
\column{.49\textwidth}
\begin{block}{Completed: exact mathematics}
Unrestricted local point classification; nine source-case audits; corrected pullbacks; regular first integrals and special branches.
\end{block}
\column{.49\textwidth}
\begin{block}{Separate: computational evidence}
Constructed local IVPs, Taylor recurrences and finite solver diagnostics. \alert{No learned integral or ML speedup yet.}
\end{block}
\end{columns}
\scope{A symmetry quotient removes independent variables ($2\to1$); a subsequent first integral lowers ODE order ($3\to2$).}
\end{frame}
\begin{frame}{Printed scalar model and regular chart}
\small
The independent variables are time $t$ and material label $s$; $\phi(t,s)$ is
particle position.  Put $p=\phi_t$ and $q=\phi_s$.  The source's two currents are
\begin{equation}
 \mathcal A=q\phi_{tt}-g\frac{\phi_{ss}}{q^2},\qquad
 \mathcal J=m\kappa\frac{p^{m-1}\phi_{ts}}{q},\qquad
 \boxed{E_m:=D_t\mathcal A-gD_s\mathcal J=0}.
 \tag{M}
\end{equation}
They follow by commuting the mixed derivatives of the source's printed bed
relations $\mathfrak b_s=-\mathcal A/g$ and $\mathfrak b_t=-\mathcal J$.
Throughout the regular calculation,
\begin{equation*}
g,\kappa>0,\qquad m\ge1,\qquad p>0,\qquad q>0.
\end{equation*}
The condition $q>0$ makes the material map locally invertible; $p>0$ fixes a
real branch of $p^{m-1}$ for noninteger $m$.
\scope{Equation (M) is the printed scalar target.  It is not the advectively corrected scalar equation or the full bed-field SVE system.}
\end{frame}
\begin{frame}[shrink=10]{One equation is not the complete physical SVE model}
\footnotesize
The manuscript audits a \textbf{printed scalar compatibility equation} in $(t,s)$, with $s$ a material label. It is \emph{not} a $(2+1)$-dimensional spatial shallow-water system.
\begin{align*}
E_m[\phi]&=D_t\mathcal A-gD_s\mathcal J,\\
\mathcal A&=\phi_s\phi_{tt}-g\frac{\phi_{ss}}{\phi_s^2},\quad
\mathcal J=m\kappa\frac{\phi_t^{m-1}\phi_{ts}}{\phi_s}.
\end{align*}
For an Eulerian bed $B(t,x)$ pulled back as $b(t,s)=B(t,\phi(t,s))$, the chain rule gives
\[
b_t= B_t+\phi_tB_x,\qquad b_s=\phi_sB_x.
\]
Thus the advective bed term cannot be discarded without an independent modelling justification. The manuscript's \emph{advectively corrected scalar compatibility} is
\[
\boxed{E_m^{\mathrm{phys}}=E_m-D_s\!\left(\frac{\phi_t\mathcal A}{\phi_s}\right)=0.}
\]
\scope{All source-case classification and ODE claims in this talk apply to the printed $E_m$, not automatically to $E_m^{\mathrm{phys}}$ or the full coupled bed-field system. [Manuscript \S7.3]}
\end{frame}
\begin{frame}[shrink=10]{First principles: the Lie group, its field, and its graphs}
\footnotesize
A \textbf{local one-parameter point group} $\{\Phi_\varepsilon\}$ satisfies $\Phi_0=\mathrm{id}$ and $\Phi_{\varepsilon+\delta}=\Phi_\varepsilon\circ\Phi_\delta$ wherever the compositions are defined. Its tangent at the identity is
\[
X=\left.\frac{d}{d\varepsilon}\right|_{0}\Phi_\varepsilon
=\tau(t,s,\phi)\partial_t+\xi(t,s,\phi)\partial_s+\eta(t,s,\phi)\partial_\phi.
\]
For a graph $\phi=f(t,s)$, transformed graph values at fixed base coordinates satisfy
\[
\boxed{\delta f=\eta-\tau f_t-\xi f_s=:Q[f].}
\]
\begin{columns}[T,onlytextwidth]
\column{.49\textwidth}
\begin{block}{Equation symmetry}
$\operatorname{pr}^{(3)}X(E_m)|_{E_m=0}=0$: the flow preserves the PDE's equation manifold and maps local solutions to solutions.
\end{block}
\column{.49\textwidth}
\begin{block}{Invariant solution}
$Q[f]=0$: this \emph{particular} graph is fixed. It selects a solution subclass; it is not the equation-symmetry test.
\end{block}
\end{columns}
\scope{$\varepsilon$ is the group-flow displacement; $m,g,\kappa$ remain fixed. A Lie \emph{algebra} consists of infinitesimal generators closed under brackets.}
\end{frame}
\begin{frame}[shrink=10]{Prolongation and free-jet splitting are the proof gate}
\footnotesize
Because $E_m$ contains third derivatives, prolong the \emph{unrestricted} $X$ through order three:
\[
\eta^t=D_t\eta-pD_t\tau-qD_t\xi,\quad
\eta^s=D_s\eta-pD_s\tau-qD_s\xi,\quad
\eta^{Ii}=D_i\eta^I-\phi_{It}D_i\tau-\phi_{Is}D_i\xi.
\]
On $p=\phi_t>0$, $q=\phi_s>0$ the leading jet $r=\phi_{ttt}$ has coefficient $q$, so the exact on-shell identity is
\[
\Omega=\left.\operatorname{pr}^{(3)}X(E_m)\right|_{r=\mathcal N}=0
\quad\text{for every remaining free jet}.
\]
\begin{center}\renewcommand{\arraystretch}{1.13}
\begin{tabular}{@{}p{1.45cm}p{4.45cm}p{4.45cm}@{}}
\toprule
Split & Consequence & Logical reason\\\midrule
$[y]\Omega=0$ & $\xi_t+p\xi_\phi=0\Rightarrow\xi=\xi(s)$ & Identity for all $p>0$\\
$[w]\Omega=0$ & $\tau_s+q\tau_\phi=0\Rightarrow\tau=\tau(t)$ & Identity for all $q>0$\\
$[ab],[a]$ & $\eta=A(t)\phi+B(t)$ & Affine form is \emph{derived}, not guessed\\
$[x]$ & $(m-3)(A-\tau')=0$, $(m-1)B'=0$ & Split $m=1,3$ \emph{before} division\\\bottomrule
\end{tabular}\end{center}
\scope{The remaining determining equations and full sufficiency identities appear in the backup. Standard infinitesimal criterion: Olver.}
\end{frame}
\begin{frame}[shrink=28]{Parameter strata: two different exceptional exponents}
\footnotesize
The complete local smooth \textbf{point} classification of the printed equation on $p,q>0$ yields:
\begin{center}\renewcommand{\arraystretch}{1.23}
\begin{tabular}{@{}lll@{}}
\toprule
Exponent & Admitted basis & Dimension\\\midrule
$m\ne1,3$ & $X_1,X_2,X_3,X_4$ & $4$\\
$m=3$ & $X_1,X_2,X_3,X_4,X_5$ & $5$\\
$m=1$ & $X_1,X_2,X_3,X_4,X_6$ & $5$\\\bottomrule
\end{tabular}\end{center}
\vspace{0.07cm}
\[
\begin{gathered}
X_1=\partial_t,\ X_2=\partial_s,\ X_3=\partial_\phi,\quad
X_4=t\partial_t+s\partial_s+\phi\partial_\phi,\\
\underbrace{X_5=t\partial_t-2s\partial_s}_{m=3\ \mathrm{only}},\qquad
\underbrace{X_6=t\partial_\phi}_{m=1\ \mathrm{only}}.
\end{gathered}
\]
\begin{block}{Why ``exceptional''?}
The rank of the determining system changes at an exponent. Neither $m=3$ symmetry nor the $m=1$ boost persists automatically at neighboring $m$.
\end{block}
\scope{The $m=1$ boost is omitted in the cited source. ``Local point algebra'' is not the algebra of the physically corrected/full SVE models. [Manuscript \S8.3]}
\end{frame}
\begin{frame}[shrink=12]{Quotient geometry: why one PDE variable disappears}
\footnotesize
Choose a verified projectable subgroup $H$. On a constant-rank regular chart require
\[
\ell=\operatorname{rank}(\text{base orbit})
=\operatorname{rank}(\text{total orbit})>0,
\qquad d=p_{\mathrm{base}}-\ell.
\]
For $X_4=t\partial_t+s\partial_s+\phi\partial_\phi$ with $t>0$, an explicit \textbf{adapted coordinate chart} is
\[
\boxed{y=\log(t/t_*),\quad z=s/t,\quad W=\phi/t},\qquad
\det\frac{\partial(y,z)}{\partial(t,s)}=t^{-2}\ne0,\quad W_\phi=t^{-1}\ne0.
\]
\begin{center}
\fcolorbox{SVEBlue}{SVEWash}{\parbox{.91\linewidth}{\centering
Orbit coordinate $y$ is forgotten by the local quotient:\\[3pt]
$\underbrace{(t,s,\phi)}_{\text{original field chart}}
\longmapsto(y,z,W)\longmapsto(z,W)$\\[4pt]
$Q=0\ \Longleftrightarrow\ \partial_yW=0\ \Longleftrightarrow\ W=F(z)$\\[3pt]
$\boxed{\phi(t,s)=tF(s/t)}$}}
\end{center}
\scope{$p_{\mathrm{base}}=2$, $\ell=1$, hence $d=1$. This describes only $H$-invariant graphs, not the whole PDE solution set; it is a \emph{local} quotient, not a global geometric quotient.}
\end{frame}
\begin{frame}[shrink=30]{Boundary conditions can invalidate a usable reduction}
\small
\textbf{Three separate gates:} equation symmetry, existence of a regular invariant graph, and compatibility with the \emph{posed data}.
\vspace{0.16cm}
For a prescribed boundary hypersurface $\Gamma=\{\rho(t,s)=0\}$, a projectable symmetry preserving that boundary must satisfy
\[
\boxed{X\rho\big|_{\Gamma}=0}.
\]
The boundary relation $\mathcal B(t,s,\phi^{(k)})=0$ must also be preserved by the appropriate prolonged action on the boundary-data manifold.
\begin{block}{Explicit counterexample: common scaling}
For $X_4=t\partial_t+s\partial_s+\phi\partial_\phi$ and a fixed initial-time surface $t=t_0>0$, take $\rho=t-t_0$. Then
\[
X_4\rho|_{t=t_0}=t_0\ne0.
\]
The PDE admits $X_4$, but the \emph{fixed-time boundary-value problem} does not inherit this symmetry.
\end{block}
\scope{Special data may still admit an invariant solution; this differs from symmetry of the whole fixed-data problem. Also check $p,q>0$, chart Jacobians, and all denominator guards.}
\end{frame}
\begin{frame}{Nine source cases: one-dimensional invariants}
\scriptsize
\renewcommand{\arraystretch}{1.16}
\begin{tabularx}{\textwidth}{@{}c p{3.0cm} p{3.2cm} X@{}}
\toprule
Case & Generator & Invariants $(z,W)$ & Invariant graph\\
\midrule
3.1 & $X_1$ & $(s,\phi)$ & $\phi=F(s)$\\
3.2 & $X_2$ & $(t,\phi)$ & $\phi=F(t)$\\
3.3 & $X_3$ & $(t,s)$ & no invariant graph\\
3.4 & $X_4$ & $(s/t,\phi/t)$ & $\phi=tF(s/t)$\\
3.5 & $X_5$, $m=3$ & $(st^2,\phi)$ & $\phi=F(st^2)$\\
3.6 & $X_1+\alpha X_2$ & $(s-\alpha t,\phi)$ & $\phi=F(s-\alpha t)$\\
3.7 & $X_1+\alpha X_3$ & $(s,\phi-\alpha t)$ & $\phi=\alpha t+F(s)$\\
3.8 & $X_2+\alpha X_3$ & $(t,\phi-\alpha s)$ & $\phi=\alpha s+F(t)$\\
3.9 & $X_1+\alpha X_2+\alpha_3X_3$ & $(s-\alpha t,\phi-\alpha_3t)$ & $\phi=\alpha_3t+F(s-\alpha t)$\\
\bottomrule
\end{tabularx}
\vspace{0.15cm}
Rows 3.7 and 3.8 correct the source's exchanged graphs.  The nine rows are an
audited selection, not a complete optimal system for every parameter stratum.
\end{frame}
\begin{frame}{Case 3.4: common scaling has an exact first integral}
\footnotesize
Let $z=s/t$, $\phi=tF(z)$, $P=F'$, $C=F''$, $H=F'''$, and $V=F-zP$.
Since $D_tG=-(z/t)DG$ and $D_sG=(1/t)DG$,
\begin{align*}
p&=V,&q&=P,&a&=z^2C/t,&b&=-zC/t,&c&=C/t,\\
r&=-(3z^2C+z^3H)/t^2,&x&=-(C+zH)/t^2,&y&=H/t^2.
\end{align*}
The currents reduce to
\[
\mathcal A=t^{-1}(z^2P-g/P^2)C,\qquad
\mathcal J=-t^{-1}m\kappa zV^{m-1}C/P.
\]
Therefore
\begin{equation*}
E_m[tF(s/t)]=t^{-2}D(K_4C),\qquad
K_4=-z^3P+\frac{gz}{P^2}+\frac{gm\kappa zV^{m-1}}P.
\end{equation*}
The third-order equation $K_4H+(DK_4)C=0$ integrates to
\begin{equation}
\boxed{K_4F''=I_0},\qquad
\boxed{F'=P,\quad P'=I_0/K_4},\quad K_4\ne0.\tag{5}
\end{equation}
The constant $I_0=K_4(z_0,F_0,P_0)C_0$ retains the third initial datum.
\end{frame}
\begin{frame}[shrink=25]{At $m=3$: the extra symmetry requires a separate chart}
\footnotesize
Only on the cubic stratum, the verified extra generator $X_5=t\partial_t-2s\partial_s$ has flow
\[
(t,s,\phi)\mapsto(e^\varepsilon t,e^{-2\varepsilon}s,\phi).
\]
On $t>0$ its independent invariants are $z=st^2$ and $W=\phi$, hence $\phi=F(st^2)$; the adapted orbit coordinate is $y=\log(t/t_*)$.
\[
\boxed{E_3[F(st^2)]=t^{-1}\mathcal R_5[F]},\qquad
\mathcal R_5[F]=2z\,D\mathcal U_5-g\,D\mathcal V_5.
\]
The third-order reduced equation is \emph{not automatically solvable in normal form}. The coefficient of $F^{(3)}$ is
\[
L_5=-\frac{2z}{(F')^2}\bigl((12g\kappa-4)z^2(F')^3+g\bigr).
\]
\begin{block}{Branch discipline}
Require $z>0$, $F'>0$ and $L_5\ne0$ before dividing. An exact singular-normal-form power profile exists on a separate undivided branch.
\end{block}
\scope{$m=3$ is an exceptional equation parameter, not an ``extra'' group-flow parameter. The common scaling $X_4$ remains admitted for every $m\ge1$.}
\end{frame}
\begin{frame}{Cases 3.6 and 3.9: traveling profiles become quadratures}
\footnotesize
Let $z=s-\alpha t$, $\phi=\alpha_3t+F(z)$,
$P=F'$, $C=F''$, $H=F'''$, and $v=\alpha_3-\alpha P>0$.
Then
\[
p=v,\ q=P,\ a=\alpha^2C,\ b=-\alpha C,\ c=C,\quad
r=-\alpha^3H,\ x=-\alpha H,\ y=H.
\]
Direct substitution into the original currents gives
\begin{align*}
E_m&=-\alpha D\{\mathcal B(P)C\},\\
\mathcal B(P)&=\alpha^2P-\frac g{P^2}
-\frac{gm\kappa(\alpha_3-\alpha P)^{m-1}}P.
\end{align*}
For $\alpha\ne0$ and $\mathcal B\ne0$,
\begin{equation}
\boxed{\mathcal B(P)F''=I_0},\qquad
\boxed{F'=P,\quad P'=I_0/\mathcal B(P)}.\tag{7}
\end{equation}
If $\mathcal Y'(P)=\mathcal B(P)$, then
\[
\mathcal Y(P(z))-\mathcal Y(P_0)=I_0(z-z_0),\qquad
F(z)=F_0+\int_{z_0}^{z}P(\zeta)\,d\zeta.
\]
The branches $I_0=0$ and $\alpha=0$ must be retained before division.
\end{frame}
\begin{frame}[shrink=12]{The correction claim is supported by exact witnesses}
\small
\textbf{Method:} substitute each \emph{printed} reduced equation into the original master residual, not into an imported simplified formula.
\vspace{0.18cm}
\begin{block}{A single decisive rational witness: source Eq. (24)}
At $m=3$, $g=1$, $\kappa=1/30$ and
\[
(z,F,F',F'',F^{(3)})=(1,2,\tfrac12,\tfrac1{10},\tfrac{23}{250}),
\]
the source's displayed reduced equation is zero, but exact original-PDE substitution gives
\[
\boxed{E_m[\phi]=\frac{2367}{5000}\ne0}\qquad(t=s=1).
\]
\end{block}
Other exact comparisons find a missing chain-rule denominator in source (27), swapped invariant graphs in cases 3.7--3.8, and errors in other displayed reductions.
\scope{A counterexample disproves universal equality of formulas, not every special solution. Source corrections and the omitted boost support a correction study; broader solver novelty is not established.}
\end{frame}
\begin{frame}[shrink=15]{Constructed local profiles: analytical and numerical checks}
\centering
\includegraphics[width=.90\linewidth]{dissertation_assets/series_profiles.pdf}
\vspace{0.06cm}
\begin{minipage}{.91\linewidth}\footnotesize
\textbf{Independent checks:} sixth-degree Taylor recurrence, classic Fehlberg 4(5), and DOP853 saved nodes for constructed cubic scaling/traveling IVPs ($g=1$, $\kappa=1/30$). Their curves overlap at this plotting scale.
\scope{Figure extracted from the supplied annotated manuscript (its Fig. 3). Finite agreement is not a certified global error or Taylor remainder bound.}
\end{minipage}
\end{frame}
\begin{frame}[shrink=30]{ML is a proposal gate, not an achieved PDE solver}
\footnotesize
After an \emph{exactly verified} profile reduction on a regular branch, set
\[
F'=P,\qquad P'=C,\qquad C'=R(z,F,P,C).
\]
A learner may propose a \emph{nonconstant} function $I(z,F,P,C)$, but it becomes a first integral only after \textbf{exact symbolic certification}:
\[
\boxed{D_zI=I_z+PI_F+CI_P+RI_C=0}.
\]
For a genuinely order-lowering chart in this representation additionally require
\[
I_C\ne0,\quad I\text{ independent of known integrals},\quad
I=I_0\text{ compatible with data}.
\]
\begin{columns}[T,onlytextwidth]
\column{.48\textwidth}\begin{block}{What is already proved?}
Classical first integrals in cases 3.4, 3.6 and 3.9, without ML.
\end{block}
\column{.48\textwidth}\begin{block}{What remains open?}
Learn a \emph{new} certifiable integral or ODE symmetry; test incremental benefit against classical reduced solvers.
\end{block}
\end{columns}
\scope{Small collocation residual, fitting a profile, and reducing neural input dimension are not proofs of lower ODE order, convergence, or computational advantage.}
\end{frame}
\begin{frame}[shrink=12]{Conclusion and defensible dissertation claims}
\small
\begin{enumerate}
\setlength{\itemsep}{0.24cm}
\item \textbf{Classification:} On the declared positive chart, the printed scalar model has point-algebra dimensions $4,5,5$ on the generic, $m=3$, and $m=1$ strata; the last contains a missing boost.
\item \textbf{Equation corrections:} Nine selected source cases are separated into regular ODEs, identities, undefined charts, and a no-graph symmetry. Rational witnesses expose erroneous printed reductions.
\item \textbf{Independent mathematical checks:} Exact first integrals, locally justified series, guarded reconstructed fields, and finite RKF45/DOP853 checks support a reproducible correction study.
\end{enumerate}
\vspace{0.15cm}
\begin{center}\color{SVEBlue}\rule{.93\linewidth}{.8pt}\end{center}
\textbf{Next contribution to test:} Can machine-proposed, \emph{symbolically certified} first integrals lower reduced-ODE order beyond established classical controls?
\scope{Not claimed: full physical SVE equivalence, global solutions, arbitrary boundary data, learned integrals, or neural speedup. Technical derivations and all numerical limits are in the backup.}
\end{frame}
\appendix
\section{Technical backup}
\begin{frame}{Backup: notation and the role of quotients}
\footnotesize
\begin{tabularx}{\linewidth}{@{}p{2.4cm}X@{}}
\toprule
Symbol & Meaning and non-equivalence\\\midrule
$(t,s)$ & Time and one material label: a $1+1$-dimensional PDE, \emph{not} two spatial dimensions\\
$\phi(t,s)$ & Particle position; $u=\phi_t$, $h=1/\phi_s$ on $\phi_s>0$ in normalized mass coordinates\\
$g,\kappa,m$ & Fixed gravity, transport coefficient, and exponent during equation-symmetry tests\\
$\varepsilon$ & One-parameter group-flow displacement, not an exponent or derivative abbreviation\\
$p,q,a,b,c$ & Jet abbreviations: $\phi_t,\phi_s,\phi_{tt},\phi_{ts},\phi_{ss}$\\
$G,\mathfrak g,H$ & Local symmetry group, its Lie algebra, and a selected subgroup (or subalgebra)\\
$y,z,W$ & Orbit coordinate, independent invariant, and fibre invariant: $(t,s,\phi)\mapsto(y,z,W)$\\
$Q=0$ & Invariant-solution graph gate, \emph{not} the PDE equation-symmetry criterion\\
\bottomrule
\end{tabularx}
\scope{The quotient forgets the $y$ direction only on a regular transverse chart and for the selected $H$-invariant solution subclass.}
\end{frame}
\begin{frame}{Backup: this Lie group is not $SO(3)$}
\footnotesize
The point-symmetry algebra for the cubic printed SVE equation is five-dimensional. The translation ideal
\[
\mathfrak n=\operatorname{span}\{X_1,X_2,X_3\},\qquad
[\mathfrak n,\mathfrak n]=0
\]
is abelian, and $[X_4,X_5]=0$, whereas $X_4,X_5$ act diagonally on translations. Hence
\[
[\mathfrak g_3,\mathfrak g_3]=\mathfrak n,\qquad
[\mathfrak n,\mathfrak n]=0.
\]
This algebra is \textbf{solvable}, with a local semidirect-product structure $\mathbb R^3\rtimes\mathbb R^2$; it cannot be $\mathfrak{so}(3)$, the three-dimensional simple rotation algebra.
\vspace{.15cm}
\begin{block}{Why discuss adjoints anyway?}
For any Lie group, conjugation acts on its Lie algebra. Here that action classifies one-dimensional subalgebras for nonredundant \emph{local} similarity reductions. The source convention is $\operatorname{Ad}(e^{\varepsilon U})V=e^{-\varepsilon\operatorname{ad}_U}V$.
\end{block}
\scope{No physical $SO(3)$ spatial rotational symmetry is being asserted. At $m=3$, $X_5$ is an anisotropic dilation of material/time coordinates.}
\end{frame}
\begin{frame}{The dissertation argument is one connected chain}
\small
The calculation begins with the equation, and every later object must pass the
previous gate:
\begin{equation*}
\boxed{E_m=0}
\xrightarrow{\;\pr^{(3)}X(E_m)|_{E_m=0}=0\;}
\boxed{\mathfrak g_m}
\xrightarrow{\;Q=0\;}
\boxed{\phi=\mathcal T[F]}
\xrightarrow{\;E_m[\mathcal T F]=0\;}
\boxed{\mathcal R[F]=0}
\xrightarrow{\;\D I=\mu\mathcal R\;}
\boxed{I=I_0}.
\end{equation*}
\vspace{0.25cm}
\begin{tabularx}{\textwidth}{@{}>{\raggedright\arraybackslash\bfseries}p{3.2cm}X@{}}
Completed result & Local point classification of the printed scalar SVE equation on a positive chart, nine audited similarity cases, corrected reduced equations, exact first integrals, local series, and finite numerical checks.\\[2mm]
Next research gate & A machine may propose an ODE first integral or reduced-ODE symmetry, but exact differentiation must certify genuine order reduction before any usefulness claim.\\
\end{tabularx}
\scope{PDE dimension reduction $2\to1$ and ODE order reduction $3\to2$ are different results.}
\end{frame}
\begin{frame}{Total differentiation produces the master equation}
\footnotesize
Introduce
\[
a=\phi_{tt},\ b=\phi_{ts},\ c=\phi_{ss},\qquad
r=\phi_{ttt},\ w=\phi_{tts},\ x=\phi_{tss},\ y=\phi_{sss}.
\]
For expressions through second order,
\begin{align*}
D_t&=\partial_t+p\partial_\phi+a\partial_p+b\partial_q
       +r\partial_a+w\partial_b+x\partial_c+\cdots,\\
D_s&=\partial_s+q\partial_\phi+b\partial_p+c\partial_q
       +w\partial_a+x\partial_b+y\partial_c+\cdots.
\end{align*}
The three nontrivial product rules are
\begin{align*}
D_t(qa)&=ab+qr,\\
D_t(cq^{-2})&=xq^{-2}-2cbq^{-3},\\
D_s(p^{m-1}bq^{-1})&=(m-1)p^{m-2}\frac{b^2}{q}
 +p^{m-1}\frac{x}{q}-p^{m-1}\frac{bc}{q^2}.
\end{align*}
\scope{The letters $a,b,c$ are derivative coordinates, not transformation parameters.}
\end{frame}
\begin{frame}{Expanded master equation and normal form}
\small
Substitution of the preceding derivatives into $E_m=D_t\mathcal A-gD_s\mathcal J$
gives
\begin{align}
E_m={}&qr+ab-g\frac{x}{q^2}+2g\frac{cb}{q^3}\notag\\
&-gm\kappa\left((m-1)p^{m-2}\frac{b^2}{q}
 +p^{m-1}\frac{x}{q}-p^{m-1}\frac{bc}{q^2}\right)=0.\tag{1}
\end{align}
Its $r$ coefficient is $q>0$, hence the equation manifold has the regular
normal form
\begin{align*}
r=\mathcal N={}&-\frac{ab}{q}+\frac{gx}{q^3}-\frac{2gcb}{q^4}\\
&+gm\kappa\left((m-1)p^{m-2}\frac{b^2}{q^2}
 +p^{m-1}\frac{x}{q^2}-p^{m-1}\frac{bc}{q^3}\right).
\end{align*}
The remaining jets are free coordinates, so their coefficients may be split in
the determining identity.
\scope{Solving for $r$ is local to $q\ne0$; it neither covers $q=0$ nor proves global PDE well-posedness.}
\end{frame}
\begin{frame}{The characteristic comes from moving a graph}
\small
Start with an unrestricted near-identity point transformation
\begin{align*}
\widetilde t&=t+\varepsilon\tau(t,s,\phi)+O(\varepsilon^2),&
\widetilde s&=s+\varepsilon\xi(t,s,\phi)+O(\varepsilon^2),\\
\widetilde\phi&=\phi+\varepsilon\eta(t,s,\phi)+O(\varepsilon^2),&
X&=\tau\partial_t+\xi\partial_s+\eta\partial_\phi.
\end{align*}
If $\phi=f(t,s)$ is a graph and $\widetilde f=f+\varepsilon\delta f+O(\varepsilon^2)$,
then
\[
\widetilde f(t+\varepsilon\tau,s+\varepsilon\xi)
=f+\varepsilon\eta+O(\varepsilon^2).
\]
Taylor expansion at the original base point yields the complete graph-change
formula
\begin{equation}
\boxed{\delta f=\eta-\tau f_t-\xi f_s},\qquad
\boxed{Q=\eta-\tau\phi_t-\xi\phi_s}.\tag{2}
\end{equation}
Thus $Q[f]=0$ selects a graph fixed by the flow.  The distinct equation-symmetry
test concerns the equation manifold, not the fixed graph; it comes next.
\end{frame}
\begin{frame}{Prolongation supplies the equation-symmetry test}
\small
Differentiating the transformed graph relation gives
\begin{align*}
\eta^t&=D_t\eta-pD_t\tau-qD_t\xi,\qquad
\eta^s=D_s\eta-pD_s\tau-qD_s\xi,\\
\eta^{Ii}&=D_i\eta^I-\phi_{It}D_i\tau-\phi_{Is}D_i\xi,\qquad
i\in\{t,s\},\quad \eta^\varnothing=\eta.
\end{align*}
Because (1) depends only on $p,q,a,b,c,r,x$, its prolonged derivative is exactly
\begin{equation*}
\pr^{(3)}X(E_m)=\eta^tE_p+\eta^sE_q+\eta^{tt}E_a+
\eta^{ts}E_b+\eta^{ss}E_c+\eta^{ttt}E_r+\eta^{tss}E_x.
\end{equation*}
The determining equation is the free-jet identity
\begin{equation}
\boxed{\Omega:=\left.\pr^{(3)}X(E_m)\right|_{r=\mathcal N}=0}.\tag{3}
\end{equation}
\scope{The recursion, not a guessed scaling ansatz, determines how $X$ acts on all derivatives.  Full third-order coefficients appear in the backup slides.}
\end{frame}
\begin{frame}{Free-jet splitting determines the coefficient functions}
\footnotesize
The restrictions are imposed in sequence; each coefficient below uses all
previous restrictions:
\begin{align*}
[y]\Omega&=\frac{g(\kappa mp^mq+p)}{pq^2}(\xi_t+p\xi_\phi)=0
 &&\Longrightarrow\quad \xi=\xi(s),\\[1mm]
[w]\Omega&=\frac{2g(\kappa mp^mq+p)}{pq^2}(\tau_s+q\tau_\phi)=0
 &&\Longrightarrow\quad \tau=\tau(t),\\[1mm]
[ab]\Omega&=-\eta_s/q=0
 &&\Longrightarrow\quad \eta=\eta(t,\phi),\\[1mm]
[a]\Omega&=q(4p\eta_{\phi\phi}-3\tau''+4\eta_{t\phi})=0
 &&\Longrightarrow\quad \eta=A(t)\phi+B(t),\quad4A'=3\tau''.
\end{align*}
The remaining $x$ coefficient then yields the complete finite system
\begin{equation}
\boxed{
2\tau'+\xi'-3A=0,\quad4A'=3\tau'',\quad
(m-3)(A-\tau')=0,\quad(m-1)A'=0,\quad(m-1)B'=0.}
\tag{4}
\end{equation}
At $m=1$, substitution into the full prolonged equation adds
$bB''+qB'''=0$, hence $B''=0$.  The factors $m-1$ and $m-3$ must be split
before division.
\end{frame}
\begin{frame}{Classification theorem and listed Lie symmetries}
\scriptsize
\begin{block}{Theorem: local smooth point algebra on $p,q>0$}
For fixed $g,\kappa>0$ and $m\ge1$, the printed scalar equation has dimension
$4$ for $m\notin\{1,3\}$, dimension $5$ for $m=3$, and dimension $5$ for
$m=1$, with the bases listed below.
\end{block}
\vspace{-0.1cm}
\begin{tabularx}{\textwidth}{@{}c>{\raggedright\arraybackslash}p{5.3cm}p{2.0cm}X@{}}
\toprule
Field & Generator & Admission & One-parameter flow $(\widetilde t,\widetilde s,\widetilde\phi)$\\
\midrule
$X_1$ & $\partial_t$ & all $m$ & $(t+\varepsilon,s,\phi)$\\
$X_2$ & $\partial_s$ & all $m$ & $(t,s+\varepsilon,\phi)$\\
$X_3$ & $\partial_\phi$ & all $m$ & $(t,s,\phi+\varepsilon)$\\
$X_4$ & $t\partial_t+s\partial_s+\phi\partial_\phi$ & all $m$ & $(e^\varepsilon t,e^\varepsilon s,e^\varepsilon\phi)$\\
$X_5$ & $t\partial_t-2s\partial_s$ & $m=3$ & $(e^\varepsilon t,e^{-2\varepsilon}s,\phi)$\\
$X_6$ & $t\partial_\phi$ & $m=1$ & $(t,s,\phi+\varepsilon t)$\\
\bottomrule
\end{tabularx}
\vspace{0.12cm}
Sufficiency follows from
$\pr X_{1,2,3}(E_m)=0$, $\pr X_4(E_m)=-2E_m$,
$\pr X_5(E_m)+E_m=(3-m)E_{\rm bed}$, and $\pr X_6(E_1)=0$.
\scope{The source omits the $m=1$ boost $X_6$.  This theorem remains local to point transformations on the stated chart.}
\end{frame}
\begin{frame}{Commutators fix the adjoint action}
\scriptsize
With $[U,V]^j=U^i\partial_iV^j-V^i\partial_iU^j$, for example
$[X_1,X_4]=X_1$ and $[X_2,X_5]=-2X_2$.  At $m=3$:
\begin{center}
\renewcommand{\arraystretch}{1.12}
\begin{tabular}{@{}c|ccccc@{}}
\toprule
$[\text{row},\text{column}]$ & $X_1$ & $X_2$ & $X_3$ & $X_4$ & $X_5$\\
\midrule
$X_1$ & $0$ & $0$ & $0$ & $X_1$ & $X_1$\\
$X_2$ & $0$ & $0$ & $0$ & $X_2$ & $-2X_2$\\
$X_3$ & $0$ & $0$ & $0$ & $X_3$ & $0$\\
$X_4$ & $-X_1$ & $-X_2$ & $-X_3$ & $0$ & $0$\\
$X_5$ & $-X_1$ & $2X_2$ & $0$ & $0$ & $0$\\
\bottomrule
\end{tabular}
\end{center}
The source convention for inner adjoints is
\[
\operatorname{Ad}(e^{\varepsilon U})V
=e^{-\varepsilon\operatorname{ad}_U}V
=V-\varepsilon[U,V]+\frac{\varepsilon^2}{2}[U,[U,V]]-\cdots.
\]
Hence $\operatorname{Ad}(e^{\varepsilon X_1})X_4=X_4-\varepsilon X_1$,
which removes a translation component.  At $m=1$, replace $X_5$ by $X_6$
and add $[X_1,X_6]=X_3$.
\scope{The two reverse-order signs in the last row correct source Table 1.}
\end{frame}
\begin{frame}{One-dimensional optimal systems}
\scriptsize
Lines are equivalent under connected inner adjoints and nonzero real line
rescaling.  For generic $m\ne1,3$, the disjoint representatives are
\begin{center}
\begin{tabular}{@{}ll@{}}
\toprule
Coefficient stratum & Representative\\
\midrule
$h\ne0$ & $\langle X_4\rangle$\\
$h=0,\ v_1\ne0$ & $\langle X_1+\nu_2X_2+\nu_3X_3\rangle$\\
$h=v_1=0,\ v_2\ne0$ & $\langle X_2+\nu_3X_3\rangle$\\
$h=v_1=v_2=0$ & $\langle X_3\rangle$\\
\bottomrule
\end{tabular}
\end{center}
At $m=3$, a nonzero dilation direction $d_4X_4+d_5X_5$ is projective; the
three resonant survivors are
\[
X_5+cX_3,\qquad X_4-X_5+cX_1,\qquad X_4+\tfrac12X_5+cX_2.
\]
Pure translations also retain support, signs, and, when $v_3\ne0$, the invariant
$K=v_1^2v_2/v_3^3$.  At $m=1$ the families are
\[
X_4+\nu X_6;\quad X_6,\ X_6+\sigma X_1+\nu X_2,\ X_6+\sigma X_2;\quad
X_1+\nu X_2,\ X_2+\nu X_3,\ X_3.
\]
\scope{This is an algebraic orbit classification.  Graph rank, equation domain, data compatibility, and solver value remain separate tests.}
\end{frame}
\begin{frame}[shrink=3]{Lemma: an invariant graph compresses the PDE}
\footnotesize
\begin{block}{Invariant-graph lemma}
Let $X$ be an admitted projectable point generator on a constant-rank chart,
with equal positive base and total orbit ranks, locally transverse
invariant graphs, and invertible base/fibre chart Jacobians. If invariants satisfy
$X(z)=X(W)=0$, then the invariant graph condition $Q=0$ integrates locally to
$W=F(z)$.  Substitution restricts $\phi(t,s)$ to one profile $F$ of one
variable; only the subsequent pullback of $E_m$ determines the profile ODE.
\end{block}
For $X_4=t\partial_t+s\partial_s+\phi\partial_\phi$,
\[
z=\frac{s}{t},\qquad W=\frac{\phi}{t},\qquad
\boxed{\phi=tF(s/t)},\qquad Q=\phi-t\phi_t-s\phi_s=0.
\]
For $X_5$ at $m=3$, $z=st^2$, $W=\phi$, hence $\phi=F(st^2)$.  For
$X_1+\alpha X_2+\alpha_3X_3$,
\[
z=s-\alpha t,\qquad W=\phi-\alpha_3t,
\qquad \boxed{\phi=\alpha_3t+F(s-\alpha t)}.
\]
\scope{The symmetry condition selects a candidate graph.  It does not itself solve the master equation.}
\end{frame}
\begin{frame}{Nine source cases: pullbacks and outcomes}
\tiny
\setlength{\tabcolsep}{4pt}
Write $P=F'$, $C=F''$, $H=F'''$, $V=F-zP$, and
\begin{align*}
K_4&=-z^3P+\frac{gz}{P^2}+\frac{gm\kappa zV^{m-1}}{P},\\
\mathcal B(P)&=\alpha^2P-\frac g{P^2}
-\frac{gm\kappa(\alpha_3-\alpha P)^{m-1}}P.
\end{align*}
\renewcommand{\arraystretch}{1.12}
\begin{tabularx}{\dimexpr\linewidth-2pt\relax}{@{}c p{4.4cm} X@{}}
\toprule
Case & Original-PDE pullback & Equation or exceptional outcome\\
\midrule
3.1 & $0$ in current form & stationary family; $P>0$, but $p=0$\\
3.2 & undefined & $q=0$ violates the master-equation domain\\
3.3 & no graph & $Q=1$ cannot vanish\\
3.4 & $t^{-2}D(K_4C)$ & $K_4H+(DK_4)C=0$; $P,V>0$\\
3.5 & $-2z\mathcal C_5/(tP^3)$ & $\mathcal C_5=0$; $m=3$, $z,P>0$\\
3.6 & $-\alpha D(\mathcal BC)$, $\alpha_3=0$ & $\mathcal BH+\mathcal B'C^2=0$ for $\alpha\ne0$\\
3.7 & $0$ & arbitrary $F(s)$; $\alpha,P>0$\\
3.8 & $\alpha H$ & $F'''=0$; $\alpha,F'>0$\\
3.9 & $-\alpha D(\mathcal BC)$ & same traveling ODE for $\alpha\ne0$; identity for $\alpha=0$\\
\bottomrule
\end{tabularx}
\vspace{0.1cm}
Regular evolution additionally requires $K_4\ne0$ in 3.4, the cubic leading
coefficient nonzero in 3.5, and $\mathcal B\ne0$ in the traveling rows.
\end{frame}
\begin{frame}{Exceptional rows have mathematically different meanings}
\small
The five nonstandard outcomes cannot be forced into one numerical template:
\begin{align*}
\text{3.1: }&\phi=F(s),\quad p=0,\quad E_m=0
 &&\text{stationary current-form branch outside }p>0,\\
\text{3.2: }&\phi=F(t),\quad q=0
 &&\text{currents undefined},\\
\text{3.3: }&X_3=\partial_\phi,\quad Q=1
 &&\text{valid equation symmetry, no fixed graph},\\
\text{3.7: }&\phi=\alpha t+F(s),\quad E_m=0
 &&\text{arbitrary smooth }F\text{ with }\alpha,F'>0,\\
\text{3.8: }&\phi=\alpha s+F(t),\quad E_m=\alpha F'''
 &&\text{exact }F=c_0+c_1t+\tfrac12c_2t^2.
\end{align*}
For example, $\phi=t+s+s^3/6$ solves 3.7 on $s>0$ although its spatial
profile has $F'''=1$.  This directly shows why exchanging rows 3.7 and 3.8
changes the mathematics.
\scope{Identity, undefined chart, no invariant graph, and regular ODE are distinct statuses, not degrees of numerical failure.}
\end{frame}
\begin{frame}{Case 3.5: cubic scaling reduces again, with chart guards}
\scriptsize
At $m=3$, use $z=st^2$ and $\phi=F(z)$.  With $\chi=12g\kappa-4$ the
undivided profile equation is
\begin{align*}
0=\mathcal C_5={}&P(\chi z^2P^3+g)H+(\chi z^2P^3-2g)C^2\\
&+(60g\kappa-12)zP^4C+(24g\kappa-2)P^5.
\end{align*}
It contains $P,C,H$ but not $F$.  Set $\mu=F'$.  The inherited $X_4$ action
on $(z,\mu)$ is $3z\partial_z-2\mu\partial_\mu$, with invariants
\[
\Sigma=z^{2/3}\mu,\qquad M=z^{5/3}\mu'.
\]
Define
\[
N(\Sigma,M)=(\chi\Sigma^3-2g)M^2
+(60g\kappa-12)\Sigma^4M+(24g\kappa-2)\Sigma^5.
\]
The corrected first-order equation is
\begin{equation}
\boxed{\frac{dM}{d\Sigma}=
\frac{\frac53M-N(\Sigma,M)/[\Sigma(\chi\Sigma^3+g)]}
     {\frac23\Sigma+M}}.\tag{6}
\end{equation}
It requires $\Sigma>0$, $\chi\Sigma^3+g\ne0$, and
$\frac23\Sigma+M\ne0$.  Source (27) omits the last chain-rule denominator.
\end{frame}
\begin{frame}{The source audit finds exact counterexamples}
\tiny
Unless stated otherwise, take $m=3$, $g=1$, $\kappa=1/30$, $z=1$,
$P=1/2$, $C=1/10$ on the positive chart.
\renewcommand{\arraystretch}{1.12}
\begin{tabularx}{\textwidth}{@{}p{1.1cm}X@{}}
\toprule
Source & Exact comparison with the original master equation\\
\midrule
(24) & Printed zero at $F=2,H=23/250$ gives $E_m=2367/5000\ne0$.\\
(25),(26) & Printed zero at $H=-88/25$ gives $E_m=437/25$; the rewrite inherits the failure.\\
(27) & Printed $M_\Sigma=89/117$; the chain rule gives $-187/65$.\\
(28),(29) & At $\alpha=-1,H=0$, printed zero gives $E_m=169/1000$.\\
(30) & Printed $M_\Xi/M=172/35$; the first integral requires $338/71$.\\
(31) & Its nonzero branch gives $(\mathcal B(P)M(P))'=960/8281\ne0$.\\
(32) & Its quadratic is sufficient but not necessary; the source uses the wrong graph.\\
(33) & At $\alpha=-1,\alpha_3=1,H=-5/119$, printed zero gives $40337/119000$.\\
(34) & The literal integral implies a jet with $E_m=-118420/328509\ne0$.\\
\bottomrule
\end{tabularx}
\vspace{0.1cm}
The exact power $F=a_\star z^{1/3}$ with $a_\star^3=45$ still solves the
undivided 3.4 and 3.5 equations, but lies where the divided 3.4 leading
coefficient vanishes.
\scope{These witnesses refute universal displayed reductions, not every special solution and not the independently verified generator algebra.}
\end{frame}
\begin{frame}{Lemma: regular profile jets have analytic series}
\footnotesize
\setlength{\abovedisplayskip}{3pt}
\setlength{\belowdisplayskip}{3pt}
\begin{block}{Local analytic existence and recurrence}
Suppose a corrected pullback has the regular form
$L(z,F,P,C)H+N(z,F,P,C)=0$ on the recorded positive branch, with
$L(z_0,F_0,P_0,C_0)\ne0$.  Then
$(F,P,C)'=(P,C,R)$, $R=-N/L$, has a unique local analytic solution germ.
For $\xi=z-z_0$ and $F=\sum_{j\ge0}a_j\xi^j$,
\end{block}
\begin{equation}
\boxed{a_{n+3}=\frac{[\xi^n]R(z_0+\xi,F_{n+2},F'_{n+2},F''_{n+2})}
{(n+3)(n+2)(n+1)}}.\tag{8}
\end{equation}
Here $a_0=F_0$, $a_1=P_0$, $a_2=C_0/2$, so
\[
a_3=R_0/6,\qquad
a_4=\frac1{24}(R_z+P_0R_F+C_0R_P+R_0R_C)_0.
\]
The degree-$N$ polynomial checks $\xi^0,\ldots,\xi^{N-3}$ exactly.  The lemma
proves only a local germ, not a global solution, convergence radius, or
remainder bound.
\end{frame}
\begin{frame}{Five regular cubic families have explicit initial series}
\footnotesize
\setlength{\abovedisplayskip}{3pt}
\setlength{\belowdisplayskip}{3pt}
Use $g=1$, $\kappa=1/30$.  Except for 3.8, the initial jet is
$(1,2,1/2,1/10)$; for 3.8 it is $(0,0,1,1/10)$.  Every row starts
\[
F(z_0+\xi)=F_0+P_0\xi+\frac{C_0}{2}\xi^2+a_3\xi^3+O(\xi^4).
\]
{\centering
\renewcommand{\arraystretch}{1.04}
\begin{tabular}{@{}ccll@{}}
\toprule
Case & Leading coefficient $L_0$ & $a_3$ & Status\\
\midrule
3.4 & $79/20$ & $-11/2370$ & common scaling\\
3.5 & $-22/5$ & $83/1100$ & cubic scaling\\
3.6 & $-71/20$ & $169/21300$ & traveling, $\alpha_3=0$\\
3.8 & $1$ & $0$ & exact $F(t)=t+t^2/20$\\
3.9 & $-79/20$ & $173/23700$ & traveling, $\alpha_3=1$\\
\bottomrule
\end{tabular}
\par}
For example,
\[
F(1+\xi)=2+\frac{\xi}{2}+\frac{\xi^2}{20}
-\frac{11\xi^3}{2370}+O(\xi^4),\qquad \phi=tF(s/t).
\]
All 21 degree-six series pass exact residual checks through $\xi^3$.
\end{frame}
\begin{frame}{Classic Fehlberg 4(5) supplies the adaptive solver}
\footnotesize
\setlength{\abovedisplayskip}{3pt}
\setlength{\belowdisplayskip}{3pt}
For $\bm y=(F,P,C)$ and $\bm y'=\bm f(z,\bm y)=(P,C,R)$, one trial step uses
six shared stages:
\begin{align*}
\bm k_i&=\bm f\!\left(z_n+c_i\Delta z,
 \bm y_n+\Delta z\sum_{j<i}a_{ij}\bm k_j\right),\quad i=1,\ldots,6,\\
\bm y_{n+1}^{(4)}&=\bm y_n+\Delta z\sum_i b_i^{(4)}\bm k_i,\qquad
\bm y_{n+1}^{(5)}=\bm y_n+\Delta z\sum_i b_i^{(5)}\bm k_i.
\end{align*}
The normalized embedded error is
\begin{equation}
e=\left[\frac13\sum_{j=1}^{3}
\left(\frac{y_{n+1,j}^{(5)}-y_{n+1,j}^{(4)}}
{\mathrm{atol}+\mathrm{rtol}\max(|y_{n,j}|,|y_{n+1,j}^{(5)}|)}\right)^2
\right]^{1/2}.\tag{9}
\end{equation}
Accept the fifth-order state when $e\le1$ and scale the next step by
$\min(5,\max(0.2,0.9e^{-1/5}))$.  Defaults are
$\mathrm{rtol}=10^{-9}$ and $\mathrm{atol}=10^{-11}$.
\scope{Stages respect all guards; SciPy RK45 is Dormand--Prince, not Fehlberg.}
\end{frame}
\begin{frame}{Convergence and reconstructed-PDE checks are finite evidence}
\footnotesize
\begin{tabularx}{\textwidth}{@{}>{\bfseries}p{3.2cm}X@{}}
\toprule
Check & Measured result and interpretation\\
\midrule
Fixed-step RKF45 & $N=1,2,4,8,16,32,64$; 55 of 147 errors resolved.  Last resolved orders are $4.97$--$4.98$ for 3.4 and $4.93$ for 3.5.  Each short traveling row supplies only one resolved pair.\\[1mm]
Tolerance sweep & Scaling endpoint errors decrease from about $10^{-7}$ to $10^{-11}$ after the step history changes.  Short traveling cases plateau because the solver retains the same two steps.\\[1mm]
Original PDE stencil & At $h=0.003$, reconstructed residual maxima are $7.98\times10^{-7}$ for DOP853 and $9.36\times10^{-7}$ for tight RKF45.  A $0.001t^3$ perturbation changes the residual by at least $0.00241$.\\
\bottomrule
\end{tabularx}
\vspace{0cm}
All 21 regular IVPs reached their endpoints under the declared sampled guards.
Finer finite-difference probes became roundoff-limited and sometimes worsened.
\scope{No finite sample proves positivity between nodes, a continuous residual bound, global continuation, or arbitrary initial-boundary compatibility.}
\end{frame}
\begin{frame}{Reconstruction gives position, velocity, and depth}
\footnotesize
The two-dimensional display uses one time coordinate and one material
coordinate, not two spatial dimensions.  On a normalized mass coordinate,
\[
x_{\rm phys}=\phi(t,s),\qquad u(t,s)=\phi_t(t,s),\qquad
h(t,s)=\frac1{\phi_s(t,s)}.
\]
\begin{center}
\renewcommand{\arraystretch}{1.04}
\begin{tabular}{@{}clll@{}}
\toprule
Case & $\phi(t,s)$ & $u(t,s)$ & $h(t,s)$\\
\midrule
3.4 & $tF(s/t)$ & $F-zF'$ & $1/F'$\\
3.5 & $F(st^2)$ & $2zF'/t$ & $1/(t^2F')$\\
3.6, 3.9 & $\alpha_3t+F(s-\alpha t)$ & $\alpha_3-\alpha F'$ & $1/F'$\\
3.7 & $\alpha t+F(s)$ & $\alpha$ & $1/F'$\\
3.8 & $\alpha s+F(t)$ & $F'$ & $1/\alpha$\\
\bottomrule
\end{tabular}
\end{center}
\scope{Only points whose invariant coordinate lies in the solved interval may be reconstructed.  These are kinematic fields; the bed equations require a separate check.}
\end{frame}
\begin{frame}{A variational Lagrangian would add a new route, not rewrite this one}
\footnotesize
Material or ``Lagrange'' coordinates and a variational Lagrangian are different
concepts.  A candidate density $\mathcal L$ would have to satisfy
\begin{equation}
\frac{\delta\mathcal L}{\delta\phi}
=\sum_{i,j\ge0}(-D_t)^i(-D_s)^j
\frac{\partial\mathcal L}{\partial\phi_{t^is^j}}
=0,
\tag{10}
\end{equation}
possibly after a separately proved equivalent reformulation of the scalar
equation.  The source reports that it could not find such a formulation; this
is not a nonexistence theorem.
\vspace{0.1cm}
If a valid action were obtained, one would still need to prove that a point
symmetry preserves $\mathcal L$ up to a total divergence.  Only then would
Noether's theorem give a conservation current, which might provide additional
first integrals or boundary diagnostics.
\vspace{0.1cm}
The reductions (5)--(7) need no action: their currents and first integrals were
checked directly against the printed master equation.
\scope{A future variational structure could strengthen conservation analysis, but it would not retroactively validate incorrect source reductions.}
\end{frame}
\begin{frame}{Replication findings and precise nonclaims}
\small
\begin{tabularx}{\textwidth}{@{}>{\bfseries}p{2.5cm}X@{}}
\toprule
Status & Statement\\
\midrule
\Exact & The unrestricted determining calculation gives dimensions $4,5,5$ on the generic, cubic, and linear-flux strata and finds the omitted $m=1$ boost.\\[1mm]
\Exact & All nine listed reductions have explicit graph and equation statuses.  Corrected 3.4, 3.6, and 3.9 equations possess verified first integrals; 3.8 is quadratic; a cubic power profile occupies a singular undivided branch.\\[1mm]
\Numerical & Twenty-one constructed regular IVPs, degree-six series, two solvers, integral checks, and reconstructed finite-difference residuals agree within the reported finite tests.\\[1mm]
Not claimed & Maximal generalized-symmetry classification, global continuation, arbitrary-data well-posedness, full bed-field equivalence, a variational action, a certified continuous numerical error, or machine-learning advantage.\\
\bottomrule
\end{tabularx}
\source{Source audit: Paliathanasis, Leon and Leach, Mathematics 14 (2026), 433.  Presentation sequence adapted from Bruz\'on and Gandarias (2021), not their transformations.}
\end{frame}
\begin{frame}{The next stage tests machine-assisted ODE order reduction}
\footnotesize
After a verified quotient, write
\[
F'=P,\qquad P'=C,\qquad C'=R(z,F,P,C).
\]
A learned proposal $I(z,F,P,C)$ is accepted as a first integral only if
\begin{equation}
\boxed{D_zI=I_z+PI_F+CI_P+RI_C=0},\qquad
\boxed{I_C\ne0},\tag{11}
\end{equation}
with $I$ nonconstant, independent of known integrals, and compatible with the
initial data.  These conditions make $I=I_0$ locally solvable for $C$ and lower
the order.  A small sampled defect or a fitted profile is insufficient.
\vspace{0.08cm}
The three research interests organize that experiment.  \textbf{(1)} The
Lagrange-coordinate option requires an explicit model-equivalence check.
\textbf{(2)} Parameter and projected-orbit-rank strata record where a symmetry
appears, disappears, or becomes unusable.  \textbf{(3)} Task-dependent
trivial-symmetry filtration applies equation, graph, data, and solver-value
gates.
\scope{Hybrid architecture design and conservation analogies are later work.  Existing first integrals are classical verified controls, not learned discoveries.}
\end{frame}
\begin{frame}{Conclusion}
\large
The replication survives because it produces three defensible results:
\vspace{0.35cm}
\begin{enumerate}
\item An unrestricted point calculation corrects the parameter-stratum
classification by adding $t\partial_\phi$ at $m=1$.
\item Direct pullback of the original master equation separates five regular
profile families from identity, undefined, and no-graph cases, and it exposes
exact errors in source equations (24)--(34).
\item Verified first integrals, Taylor recurrences, and independent numerical
checks make the corrected local ODEs reproducible and provide the controls for
future machine-assisted order reduction.
\end{enumerate}
\vfill
{\color{SVEBlue}\rule{\textwidth}{0.7pt}}\par\smallskip
{\small\color{SVEMuted}Every conclusion retains its model, parameter stratum,
chart, denominator, and exact-versus-numerical status.}
\end{frame}
\begin{frame}{Selected references}
\scriptsize
\setbeamertemplate{bibliography item}[text]
\begin{thebibliography}{99}
\setlength{\itemsep}{1pt}
\bibitem{ap} A. Paliathanasis, G. Leon, and P. G. L. Leach,
\emph{Lie Symmetries and Similarity Solutions for a Shallow-Water Model with Bed Elevation in Lagrange Variables},
Mathematics 14 (2026), 433.
\bibitem{bg} M. S. Bruz\'on and M. L. Gandarias,
\emph{Symmetry Analysis and Conservation Laws of a Family of Boussinesq Equations},
in \emph{Symmetries and Applications of Differential Equations}, Springer, 2021, pp. 153--171.
\bibitem{olver} P. J. Olver, \emph{Applications of Lie Groups to Differential Equations}, 2nd ed., Springer, 1993.
\bibitem{teschl} G. Teschl, \emph{Ordinary Differential Equations and Dynamical Systems}, AMS, 2012.
\bibitem{uddin} M. J. Uddin et al., \emph{Two parameters Lie group analysis and numerical solution of unsteady free convective flow of non-Newtonian fluid}, Alexandria Engineering Journal, 2016.
\bibitem{fn} E. D. Fern\'andez-Nieto et al., \emph{Formal deduction of the Saint-Venant--Exner model including arbitrarily sloping sediment beds and associated energy}, ESAIM: M2AN 51 (2017), 115--145.
\bibitem{noether} E. Noether, \emph{Invariant Variation Problems}, 1918; English translation, 2005.
\bibitem{otto} S. E. Otto et al., \emph{A Unified Framework to Enforce, Discover, and Promote Symmetry in Machine Learning}, JMLR 26 (2025), 1--83.
\end{thebibliography}
\end{frame}
\begin{frame}{Backup: higher prolongation coefficients}
\scriptsize
The second- and third-order coefficients actually used in (3) are
\begin{align*}
\eta^{tt}&=D_t^2\eta-pD_t^2\tau-qD_t^2\xi-2aD_t\tau-2bD_t\xi,\\
\eta^{ts}&=D_tD_s\eta-pD_tD_s\tau-qD_tD_s\xi
-aD_s\tau-bD_s\xi-bD_t\tau-cD_t\xi,\\
\eta^{ss}&=D_s^2\eta-pD_s^2\tau-qD_s^2\xi-2bD_s\tau-2cD_s\xi,\\
\eta^{ttt}&=D_t\eta^{tt}-rD_t\tau-wD_t\xi,\\
\eta^{tss}&=D_t\eta^{ss}-xD_t\tau-yD_t\xi.
\end{align*}
The necessary master partials include
\begin{align*}
E_a&=b,\qquad
E_b=a+\frac{2gc}{q^3}-\frac{2gm\kappa(m-1)p^{m-2}b}{q}
+\frac{gm\kappa p^{m-1}c}{q^2},\\
E_c&=\frac{2gb}{q^3}+\frac{gm\kappa p^{m-1}b}{q^2},\qquad
E_r=q,\qquad E_x=-\frac g{q^2}-\frac{gm\kappa p^{m-1}}q.
\end{align*}
The free coordinates $w$ and $y$ enter through these prolonged coefficients
even though the original $E_m$ has no explicit $w$ or $y$ term.
\end{frame}
\begin{frame}{Backup: classic Fehlberg tableau}
\scriptsize
\begin{center}
\renewcommand{\arraystretch}{1.18}
\begin{tabular}{@{}c|rrrrrr@{}}
\toprule
$c_i$ & \multicolumn{6}{c}{$a_{ij}$}\\
\midrule
$0$ & $0$ & & & & &\\
$1/4$ & $1/4$ & $0$ & & & &\\
$3/8$ & $3/32$ & $9/32$ & $0$ & & &\\
$12/13$ & $1932/2197$ & $-7200/2197$ & $7296/2197$ & $0$ & &\\
$1$ & $439/216$ & $-8$ & $3680/513$ & $-845/4104$ & $0$ &\\
$1/2$ & $-8/27$ & $2$ & $-3544/2565$ & $1859/4104$ & $-11/40$ & $0$\\
\midrule
$b^{(4)}$ & $25/216$ & $0$ & $1408/2565$ & $2197/4104$ & $-1/5$ & $0$\\
$b^{(5)}$ & $16/135$ & $0$ & $6656/12825$ & $28561/56430$ & $-9/50$ & $2/55$\\
\bottomrule
\end{tabular}
\end{center}
The implementation propagates the fifth-order estimate and uses the
fourth--fifth difference only as the local adaptive indicator.
\end{frame}
\begin{frame}{Backup: symbolic replication workflow}
\footnotesize
\begin{enumerate}
\item Expand $E=D_t\mathcal A-gD_s\mathcal J$ in third-order jets and solve for the permitted leading derivative.
\item Introduce unrestricted $\tau,\xi,\eta$, form $Q$, construct $\eta^I$, and split $\pr^{(3)}X(E)|_{E=0}$ in free jets.
\item Integrate separately on $m=1$, $m=3$, and $m\ne1,3$; substitute every basis field into the full prolonged equation.
\item Compute brackets and adjoints, then solve $X(z)=X(W)=0$ for each selected subgroup.
\item Differentiate each graph through order three in $(t,s)$ and substitute it into the original master equation.
\item Check every proposed first integral by exact total differentiation before dividing by any coefficient.
\item Compare with the printed formulas through exact substitutions; retain every nonzero witness.
\item At regular initial jets, construct Taylor coefficients, verify the undivided residual, solve by classic RKF45, and compare with DOP853 and lower-order equations.
\item Reconstruct $\phi,u,h$ only on the solved chart and evaluate the original two-variable residual independently.
\end{enumerate}
\scope{The program performs these symbolic operations; source formulas are comparison inputs rather than trusted executable equations.}
\end{frame}
\begin{frame}{Backup: series and numerical profiles}
\centering\includegraphics[width=.95\linewidth]{dissertation_assets/series_profiles.pdf}
\scope{Extracted from the supplied annotated manuscript, Figure 3 (top panels).}
\end{frame}
\begin{frame}{Backup: Taylor differences at finite nodes}
\centering\includegraphics[width=.92\linewidth]{dissertation_assets/series_errors.pdf}
\scope{Extracted from the supplied annotated manuscript, Figure 3 (bottom panels). Log-scale differences are not a certified continuous error estimate.}
\end{frame}
\begin{frame}{Backup: first-integral diagnostics}
\centering\includegraphics[width=.92\linewidth]{dissertation_assets/first_integral_checks.pdf}
\scope{Extracted from the supplied annotated manuscript, Figure 4. Drift and solution differences near roundoff are finite diagnostics.}
\end{frame}
\end{document}
